\documentclass[11pt]{article}
\usepackage[T1]{fontenc}
\usepackage{microtype}
\usepackage[margin=1in]{geometry}
\usepackage{amsmath,amssymb,amsthm}
\usepackage[hidelinks]{hyperref}
\newtheorem{theorem}{Theorem}
\newtheorem{lemma}[theorem]{Lemma}
\newtheorem{corollary}[theorem]{Corollary}
\newtheorem{proposition}[theorem]{Proposition}
\DeclareMathOperator{\CT}{CT}
\DeclareMathOperator{\Newt}{Newt}
\newcommand{\A}{\mathcal O}
\newcommand{\Gm}{\mathbb G_m}
\newcommand{\dd}{\mathrm d}
\newcommand{\DR}{\mathrm{dR}}
\title{Finite Tests for Mixed Moments of Regular Functions}
\author{Roy van Rijn}
\date{}
\hypersetup{pdftitle={Finite Tests for Mixed Moments of Regular Functions},pdfauthor={Roy van Rijn}}

\begin{document}
\maketitle

\begin{abstract}
Let \(X\) be a smooth affine variety in characteristic zero with a
nowhere-vanishing top form. Let \(\Lambda\) be a linear functional on
regular functions that annihilates exact forms after multiplication by
this top form. For any fixed tuple of regular functions, we give a
finite test for whether the mixed moments against every fixed regular
multiplier vanish whenever all exponents are sufficiently large.
A generating series records every offset from a common diagonal
exponent. The resulting problem uses one exponent shift in twisted
de Rham cohomology. A differential equation bounds the number of
coefficients of each series that must be checked. For Laurent constant
terms, a single common weight characterizes the vanishing property,
and every failure occurs infinitely often on some positive integer ray.
\end{abstract}

\section{A diagonal can miss the mixed moments}

Take \(f_1=x^2\) and \(f_2=x^{-1}\) on the one-dimensional torus, and
let \(\Lambda=\CT_x\) be Laurent constant term. On the diagonal,
\[
   \CT_x(h f_1^m f_2^m)=\CT_x(hx^m)=0
   \qquad\text{for all sufficiently large \(m\), for every fixed \(h\)}.
\]
Yet \(\CT_x(f_1^m f_2^{2m})=1\) for every \(m\ge1\).
We test a quadrant of exponents at once by expanding a product of
geometric series. Independent variables keep the coefficients for
different offsets distinct.

Let \(K\) be a field of characteristic zero, let \(X\) be a smooth
geometrically integral affine \(K\)-variety of dimension \(n\ge1\),
and let \(A=\A(X)\). Fix a nowhere-vanishing regular \(n\)-form
\(\omega\). Throughout, \(\Lambda:A\to K\) is \(K\)-linear and satisfies
\begin{equation}\label{eq:trace}
 \Lambda(g)=0
 \quad\text{whenever}\quad g\omega=\dd\beta,\qquad
 \beta\in\Gamma(X,\Omega_X^{n-1}).
\end{equation}
Over \(\mathbb C\) this holds for a period over a closed \(n\)-cycle
in \(X(\mathbb C)\), by Stokes' theorem. It also holds for an
invariant functional on a smooth affine
homogeneous space with invariant volume. The infinitesimal action
spans the tangent bundle, so every \((n-1)\)-form is a sum of forms
\(g\,\iota_\xi\omega\) with \(\xi\) an infinitesimal action.
Here \(\dd(g\,\iota_\xi\omega)=\xi(g)\omega\), and invariance
annihilates \(\xi(g)\). For example, for a nondegenerate quadratic
form \(q\), the residue of the ambient volume divided by \(q-1\)
gives an \(\mathrm{SO}(q)\)-invariant volume on the norm-one quadric.

For regular \(f_1,\ldots,f_s\), write \(f^{\mathbf m}=\prod_j f_j^{m_j}\)
and \(F=\prod_j f_j\). Say that the tuple has \emph{cofinal vanishing}
if for every fixed \(h\in A\) there is an integer \(M_h\) such that
\(\Lambda(hf^{\mathbf m})=0\) whenever \(m_j\ge M_h\) for every \(j\).
The cutoff can depend on \(h\).

\begin{theorem}\label{thm:main}
Assume every \(f_j\) is nonzero. There are an integer \(M\ge1\),
regular functions \(h_1,\ldots,h_r\), and integers \(b_1,\ldots,b_r\ge0\)
such that cofinal vanishing is equivalent to the finite list of tests
\begin{equation}\label{eq:finite-test}
 \Lambda(h_i F^M f^{\boldsymbol\alpha})=0
 \qquad
 (1\le i\le r,\quad
   \boldsymbol\alpha\in\mathbb Z_{\ge0}^s,\quad |\boldsymbol\alpha|\le b_i).
\end{equation}
Given \(X,\omega,f_1,\ldots,f_s\) over a finitely generated field
presented over \(\mathbb Q\), the list is computable. Deciding the
tests uses exact evaluation of \(\Lambda\) on regular functions.
\end{theorem}

A zero phase makes cofinal vanishing immediate. The proof below also
gives a finite test for cofinal vanishing with one prescribed
multiplier. For Laurent constant terms we obtain a simpler answer
in Section~\ref{sec:laurent}.

The exponent shift is the algebraic Mellin shift of Loeser and Sabbah
\cite[\S1.3]{LS}. For a single phase on quadratic complements, it gives
a test at one sufficiently large order \cite{QuadraticComplements}.
The new step here records an entire quadrant by a formal expansion
and then reduces the series test to ordinary moments.

\section{The exponent shift}\label{sec:shift}

We use the following fact over an arbitrary characteristic-zero
coefficient field \(L\). The output field of the functional may be
an extension \(C\) of \(L\).

\begin{lemma}\label{lem:shift}
Suppose \(B\) is the coordinate ring of a smooth integral affine
\(L\)-variety of dimension \(n\), with volume form \(\eta\).
Let \(\lambda:B\to C\) be \(L\)-linear and annihilate exact forms
as in \eqref{eq:trace}. For \(P\in B\) nonzero, let
\[
 H= \Omega^n_{B[P^{-1}]/L}(\nu)
       \big/\nabla_\nu\Omega^{n-1}_{B[P^{-1}]/L}(\nu),
 \qquad
 \nabla_\nu=\dd+\nu\,\dd P/P\wedge .
\]
Here \((\nu)\) denotes scalar extension to \(L(\nu)\).
The space \(H\) is finite-dimensional. It has a basis represented
by \(a_i\eta\) with \(a_i\in B\). For any such basis, the vector
\(u_m=(\lambda(a_iP^m))_i\) satisfies
\[
                     u_{m+1}=A(m)u_m
\]
with \(A(\nu)\) a rational matrix invertible at every sufficiently
large integer \(m\). For a fixed \(a\in B\), the scalar sequence
\(\lambda(aP^m)\) also satisfies a polynomial recurrence for every
sufficiently large \(m\).
\end{lemma}

\begin{proof}
Finiteness follows from the logarithmic model in
Section~\ref{sec:effective}. A form \(g(\nu)\eta\) represents the
sequence \(\lambda(P^m g(m))\) modulo eventual equality. The sequence
is defined for all sufficiently large \(m\), since there are finitely
many parameter poles and \(P^m\) clears the powers of \(P^{-1}\).
A rational scalar \(r(\nu)\) acts on such sequences by multiplication
by \(r(m)\) for large \(m\). If \(g(\nu)\eta=\nabla_\nu\beta\), then
\[
                P^m\nabla_m\beta(m)=\dd(P^m\beta(m)).
\]
The form inside the last differential is regular once \(m\) is large,
so the sequence is eventually zero. The assignment therefore factors
through \(H\).

The map \(T\gamma(\nu)=P\gamma(\nu+1)\) satisfies
\(\nabla_\nu T\gamma=T\nabla_\nu\gamma\), and its inverse is
\(P^{-1}\gamma(\nu-1)\). Thus \(T\) is a semilinear automorphism of
\(H\), and it shifts the corresponding sequences by one. Apply a
common power of \(T\) to any basis to clear all powers of \(P^{-1}\).
Expanding the resulting numerators over \(L(\nu)\) shows that classes
represented by \(a\eta\) with \(a\in B\) span. Select a basis from them.

Write \(T[a_i\eta]=\sum_j A_{ij}(\nu)[a_j\eta]\). The matrix is
invertible because \(T\) is. Specialize the finitely many relations
beyond their parameter poles and denominator cutoffs, and avoid the
integer zeros of the determinant. This gives the vector recurrence.
For a fixed \(a\), linear dependence among
\([a\eta],[Pa\eta],\ldots,[P^r a\eta]\), where \(r=\dim H\), gives
a scalar recurrence after clearing denominators.
\end{proof}

For a basis as in Lemma~\ref{lem:shift}, all the sequences
\(\lambda(aP^m)\), \(a\in B\), eventually vanish exactly when \(u_M=0\)
at any single order \(M\) beyond the matrix cutoff. Indeed the zero
vector propagates in both directions there, and the basis spans the
quotient through which the sequences factor.

\section{Recording all offsets}\label{sec:offsets}

Take indeterminates \(u_1,\ldots,u_s,t\), and put
\(E=K(u_1,\ldots,u_s)\), \(L=E(t)\), and \(G_j=1-tu_j f_j\).
Work with the algebra
\begin{equation}\label{eq:localization}
        B=(A\otimes_K L)[G_1^{-1},\ldots,G_s^{-1}].
\end{equation}
Expand at \(t=0\). Each \(G_j\) has constant coefficient \(1\),
so this gives a map to \((A\otimes_K E)((t))\). Applying \(\Lambda\)
coefficientwise gives an \(L\)-linear functional
\(\lambda_t:B\to E((t))\). Expansion commutes with \(d_X\), the
differential in the coordinates of \(X\). Consequently \(\lambda_t\)
annihilates the exact relative top forms on \(\operatorname{Spec} B\).

For \(h\in A\) and \(M\ge1\), the basic identity is
\begin{equation}\label{eq:offsets}
 \lambda_t\left(\frac{hF^M}{\prod_jG_j}\right)
 =\sum_{\boldsymbol\alpha\ge0}
    \Lambda(hF^M f^{\boldsymbol\alpha})\,
    u^{\boldsymbol\alpha}t^{|\boldsymbol\alpha|}.
\end{equation}
For each power of \(t\), distinct offsets give distinct monomials in
the independent \(u_j\). The series is zero exactly when all the
mixed moments in the quadrant \(m_j\ge M\) vanish for this \(h\).

Cofinal vanishing on \(A\) is equivalent to
\begin{equation}\label{eq:equiv}
       \lambda_t(aF^m)=0\quad\text{eventually for every fixed \(a\in B\)}.
\end{equation}
To see the forward implication, express \(a\) as a finite sum of
\(L\)-multiples of \(h/\prod_jG_j^{e_j}\) with \(h\in A\) and
\(e_j\ge0\). The coefficients of the formal expansion of this last
expression times \(F^m\) involve the mixed moments for \(h\) at
exponents at least \(m\). The cutoff for \(h\) therefore works for
every coefficient at once. Take the maximum cutoff over the finitely
many numerators. The reverse implication follows from
\eqref{eq:offsets} applied to each fixed \(h\).

Apply Lemma~\ref{lem:shift} over \(L\) with phase \(P=F\).
After clearing the \(F\)-denominators as in its proof, choose a common
denominator \(Q=\prod_j G_j^{e_j}\) for the basis representatives
with all \(e_j\ge1\). Expanding their numerators over \(L(\nu)\)
shows that some \(h_i/Q\), with \(h_i\in A\), represent a basis.
The matrix recurrence tests \eqref{eq:equiv} at one order \(M\)
for these representatives. The test series are
\begin{equation}\label{eq:weighted-offsets}
 W_i(t)=\lambda_t\left(\frac{h_i F^M}{Q}\right)
   =\sum_{\boldsymbol\alpha\ge0}
     \left(\prod_j\binom{e_j+\alpha_j-1}{\alpha_j}\right)
     \Lambda(h_i F^M f^{\boldsymbol\alpha})\,
     u^{\boldsymbol\alpha}t^{|\boldsymbol\alpha|}.
\end{equation}
The binomial factors are nonzero in characteristic zero. The
coefficient test below supplies a bound \(b_i\) for each series,
which proves Theorem~\ref{thm:main}.

\section{From formal series to ordinary moments}\label{sec:jets}

\begin{lemma}\label{lem:ode}
For a fixed \(g\in B\) such that \(W(t)=\lambda_t(g)\) belongs to
\(E[[t]]\), there is a computable integer \(b\ge0\) such that \(W=0\)
if and only if its coefficients of orders \(0,\ldots,b\) vanish.
Effectivity assumes that the algebraic data are given over a
presented finitely generated field.
\end{lemma}

\begin{proof}
The classes of \(\partial_t^j(g\omega)\) in the untwisted relative
top de Rham space of \(\operatorname{Spec} B\) are linearly dependent
over \(E(t)\). Here \(\partial_t\) fixes \(A\) and \(E\); it extends
to the localization and commutes with \(d_X\). Clear denominators
in a dependence to obtain
\[
             \sum_j p_j(t)\partial_t^j(g\omega)=d_X\beta,
             \qquad p_j\in E[t],\quad \text{some }p_j\ne0 .
\]
The coefficientwise functional and \(\partial_t\) commute, so \(W\)
satisfies the same nonzero homogeneous differential equation.
The relation can be found by increasing bounds on polynomial degrees,
derivative order, and the denominators of \(\beta\), and solving linear
systems after clearing denominators.

For each nonzero \(p_j\), let \(v_j=\operatorname{ord}_t(p_j)\)
and \(c_j=[t^{v_j}]p_j\). Put
\[
  \delta=\min_{p_j\ne0}(v_j-j),\qquad
  I(k)=\sum_{v_j-j=\delta}c_j\,k(k-1)\cdots(k-j+1).
\]
The empty product for \(j=0\) is \(1\). The degrees \(j\) in this sum
are distinct, so \(I\) is nonzero. If \(W\) has first nonzero
coefficient \(w_k\), the coefficient of \(t^{k+\delta}\) in its
differential equation is \(I(k)w_k\). This identity also holds when a
derivative kills the leading term, by the falling factorial in the
formula. Hence \(I(k)=0\). A bound for the nonnegative integer roots
of \(I\) is a valid \(b\). It can be computed from any nonzero
rational coordinate polynomial of \(I\). To obtain one, express
the coefficients in a basis over a rational function field in a
transcendence basis, clear the common denominators, and choose a
nonzero rational coefficient of the resulting monomials.
\end{proof}

This is the usual Picard--Fuchs dependence of differentiated forms
\cite{Lairez}, used here with a formal coefficientwise functional.
For \eqref{eq:weighted-offsets}, checking orders \(0,\ldots,b_i\)
is precisely checking the moments with \(|\boldsymbol\alpha|\le b_i\).

\subsection*{A worked test}
Return to \(f_1=x^2\), \(f_2=x^{-1}\). Write \(a=tu_1\),
\(b=tu_2\) and \(Q=(1-a x^2)(1-b/x)\).
Since \(F=x\), the twisted quotient in the frame \(\dd x/x\)
is the cokernel of \(x\partial_x+\nu\).
Over a splitting field, \(Q\) has three distinct nonzero roots.
Applying \(x\partial_x+\nu\) raises a pole of order \(j\) there
to order \(j+1\) with nonzero leading coefficient.
Thus higher poles reduce to simple poles, and no nonzero
simple-pole part alone is exact. Laurent polynomial terms are
exact since \(k+\nu\) is invertible for every integer \(k\).
The values of \(1,x,x^2\) at the three roots give an invertible
Vandermonde matrix, so \(1/Q,x/Q,x^2/Q\) represent a basis.

From \(xQ=x-b-a x^3+abx^2\), their sequence vector has shift matrix
\[
       \begin{pmatrix}
           0&1&0\\
           0&0&1\\
           -b/a&1/a&b
         \end{pmatrix},
       \qquad \det=-b/a.
\]
The only Laurent term used in this reduction is a multiple of \(x\),
whose exactness certificate has denominator \(\nu+1\). Thus \(M=1\)
is allowed. The first basis representative already detects failure.
Its series at order \(M\) is \(W=(tu_2)^M/(1-ct^3)\), where
\(c=u_1u_2^2\), and satisfies
\[
            t(1-ct^3)W'-(M+(3-M)ct^3)W=0.
\]
The indicial polynomial is \(k-M\), and the coefficient of order
\(M\) is \(u_2^M\). For the other representatives \(x^i/Q\) the same
formula replaces \(M\) by \(M+i\), giving bounds \(1,2,3\) for the
three series at \(M=1\). The first already detects the ordinary
moment \(\CT_x(f_1 f_2^2)=1\).

\begin{corollary}\label{cor:fixed}
For a prescribed \(h\), cofinal vanishing of
\(\Lambda(hf^{\mathbf m})\) is decidable under the same effective
hypotheses. Vanishing at every strictly positive exponent vector
is also decidable for this \(h\).
\end{corollary}

\begin{proof}
Let \(B_m\) be the series in \eqref{eq:offsets} with \(M=m\).
Lemma~\ref{lem:shift} gives a scalar recurrence for \(B_m\).
Drop unused shifts at the two ends. Beyond all specialization
exceptions and the integer zeros of the first coefficient,
eventual zero propagates backwards. If \(B_M=0\) at one order,
every \(B_m\) for \(m\ge M\) is zero because the corresponding
quadrants are nested. Thus Lemma~\ref{lem:ode} applied at any order
\(M\) past the recurrence cutoff decides cofinal vanishing.
To test every strictly positive exponent vector, apply the same
series test directly to \(B_1\).
\end{proof}

\section{Computing the cohomology data}\label{sec:effective}

Here is a finite computation for the cohomology used above.
It applies both to the logarithmic twist and to the untwisted
complex in Lemma~\ref{lem:ode}.

Let \(V\) be the smooth affine variety on which the forms are
computed. Take a smooth projective compactification \(Y\) with
simple normal crossings boundary \(D=Y\setminus V\).
Starting from the projective closure of an affine embedding,
resolve its singularities and principalize the total boundary
ideal \cite[Theorems 1.6 and 1.10]{BM}. The centers are outside
\(V\), so the blowup charts give the compactification and its
identification on \(V\). For the twist, the
function \(P\) is a unit on \(V\), so \(\dd P/P\) is logarithmic
along \(D\). On the logarithmic complex with coefficients
\(\A_Y(bD)\), for \(b\ge0\), the residue along a component \(D_i\)
is \(\nu\,\operatorname{ord}_{D_i}(P)-b\). No residue is a
strictly positive integer over the field with indeterminate
\(\nu\). The untwisted case has residues \(-b\).
The logarithmic comparison therefore computes the required
algebraic de Rham space from
\(\Omega_Y^\bullet(\log D)(bD)\) with the corresponding differential
\cite[II, Corollaires 3.14--3.15]{Deligne}. For characteristic-zero
fields this application can be checked on a finitely generated
model embedded in \(\mathbb C\) and passed back by flat base change.
The hypercohomology is finite-dimensional by projectivity and
coherence of the terms.

For effectivity, embed \(Y\) in projective space and let \(H_i\)
be the divisors of its homogeneous coordinates in a
nondegenerate projective embedding. For nonempty
subsets \(I\), use the finite Čech terms
\[
 \Gamma\!\left(Y,\,
   \Omega_Y^p(\log D)\left(bD+(a+p)\sum_{i\in I}H_i\right)\right)
       \qquad\text{in Čech degree } |I|-1.
\]
The Čech maps are signed inclusions. The differential increases
the allowed pole order along each \(H_i\) by at most one, so it
preserves the double complex. For each \(p\), the augmented sheaf
complex of these terms is exact. Locally there is an \(H_i\)
missing the point, and inserting that index gives a contracting
homotopy. For large \(a\) the non-augmented sheaf terms are acyclic
by Serre vanishing, since all the \(H_i\) are hyperplane divisors.
One can search for such an \(a\) using coherent sheaf cohomology
computed from graded presentations. The finite total complex then
computes the hypercohomology. This construction uses the differential
on rational forms explicitly.

For a regular top form on \(V\), increase \(b\) until the form
is a global section of \(\Omega_Y^n(\log D)(bD)\).
In the finite model it enters diagonally in Čech degree zero.
Its coordinates in the section spaces are found by clearing
denominators and comparing rational forms on affine charts.
Increasing \(a,b\) when necessary gives inclusion maps
preserving the computed cohomology. This permits basis selection
from localized monomials. On the affine variety, a linear relation in top cohomology
can be certified by a differential of an \((n-1)\)-form. Searching
over increasing degree and denominator bounds finds such certificates
by linear algebra. Their denominators give the specialization
cutoffs in Lemma~\ref{lem:shift}. Integer roots of determinant
numerators and parameter denominators are bounded by the rational
coordinate argument used in Lemma~\ref{lem:ode}.

\section{Laurent constant terms}\label{sec:laurent}

For a nonzero Laurent polynomial \(f\), let \(\Newt(f)\) be the
convex hull of its exponent vectors. The Laurent constant term
fits \eqref{eq:trace} on a torus with volume form
\(\bigwedge_j \dd x_j/x_j\). In this case the mixed test has
a direct polyhedral form.

\begin{theorem}\label{thm:laurent}
Let \(f_1,\ldots,f_s\) be nonzero Laurent polynomials over a
characteristic-zero field and put \(P_i=\Newt(f_i)\).
Cofinal vanishing against every fixed Laurent multiplier holds
exactly when there is an integral linear weight nonnegative on
every \(P_i\) and strictly positive on at least one entire \(P_i\).
If this condition fails, there are positive integers \(d_i\)
such that \(\CT(\prod_i f_i^{k d_i})\ne0\) for infinitely many \(k\).
\end{theorem}

\begin{proof}
With such a weight, the minimum weight of the mixed product
tends to infinity as all exponents tend to infinity. Adding
the fixed minimum weight of any multiplier still gives strictly
positive weight once the exponents are large, so its constant
term is zero.

Let \(P=\sum_i P_i\) and let \(C\) be the cone generated by the union
of the \(P_i\). The rational polyhedron \(P+C\) is closed. If it
excludes the origin, strict separation gives a rational weight
nonnegative on \(C\) and strictly positive on \(P\). Clearing
denominators gives the required integral weight. Thus failure
of the weight condition gives \(0\in P+C\). By convexity,
\[
                P+C=\bigcup_{a_i\ge1}\sum_i a_i P_i .
\]
Indeed adding cone summands from \(P_i\) amounts to enlarging its
coefficient in the Minkowski sum. More explicitly, solve
\(\sum_{i,v}\lambda_{iv}v=0\) with \(\lambda_{iv}\ge0\) and
\(\sum_v\lambda_{iv}\ge1\) for each \(i\), where \(v\) runs over
the vertices of \(P_i\). The feasible system is rational, and
\(a_i=\sum_v\lambda_{iv}\) gives the required rational coefficients.
After scaling, we obtain positive integers \(d_i\) with
\(0\in\sum_i d_iP_i=\Newt(\prod_i f_i^{d_i})\).

The theorem of Duistermaat and van der Kallen says that a Laurent
polynomial whose positive powers all have zero constant term has
Newton polytope excluding the origin \cite[\S3]{DvK}.
It applies over the coefficient field here by embedding the
finitely generated subfield of coefficients in \(\mathbb C\).
If the constant terms of powers of \(\prod_i f_i^{d_i}\) were
eventually zero, take a fixed sufficiently large power.
All positive powers of that power would have zero constant term,
while its Newton polytope still contains the origin.
This contradiction gives the claimed ray.
\end{proof}

For example \(x(1+y)\) and \(y+y^{-1}\) satisfy the criterion
by weighting the \(x\)-exponent. The second phase has weight zero
throughout its support. For the phases \(x^2,x^{-1}\) in the
opening example, the weight condition fails, and \((1,2)\)
is a witness direction.

For general traces, a structural question is whether failure of
cofinal vanishing against all multipliers always gives a single
positive integer ray on which the moments for some fixed multiplier
are nonzero infinitely often.

\paragraph{AI assistance and review.}
ChatGPT (OpenAI) assisted with mathematical arguments, literature
checks, exposition, proof review and verification code.

\begin{thebibliography}{9}
\bibitem{BM}
E. Bierstone and P.~D. Milman,
\emph{Canonical desingularization in characteristic zero by blowing up
the maximum strata of a local invariant},
Invent. Math. \textbf{128} (1997), 207--302.
\href{https://arxiv.org/abs/alg-geom/9508005}{arXiv:alg-geom/9508005}.

\bibitem{Deligne}
P. Deligne, \emph{\'Equations diff\'erentielles \`a points singuliers
r\'eguliers}, Lecture Notes in Mathematics \textbf{163},
Springer, 1970.
\href{https://publications.ias.edu/book/export/html/355}{Author's text}.

\bibitem{DvK}
J.~J. Duistermaat and W. van der Kallen,
\emph{Constant terms in powers of a Laurent polynomial},
Indag. Math. (N.S.) \textbf{9} (1998), 221--231.
\href{https://webspace.science.uu.nl/~kalle101/powers.pdf}{Author's copy}.

\bibitem{Lairez}
P. Lairez, \emph{Computing periods of rational integrals},
Math. Comp. \textbf{85} (2016), 1719--1752.
\href{https://arxiv.org/abs/1404.5069}{arXiv:1404.5069}.

\bibitem{LS}
F. Loeser and C. Sabbah,
\emph{\'Equations aux diff\'erences finies et d\'eterminants d'int\'egrales
de fonctions multiformes},
Comment. Math. Helv. \textbf{66} (1991), 458--503.
\href{https://doi.org/10.1007/BF02566659}{DOI: 10.1007/BF02566659}.

\bibitem{QuadraticComplements}
R. van Rijn, \emph{The Mathieu property for quadratic complements},
manuscript.
\end{thebibliography}
\end{document}
